%%%%%%%%%%%%%%%%%%%%%%%%%%%%%%%%%%%%%%%%
%
% $ Autor: Gruppe 07 $
% $ Projekt: Magic Faucet Fountain $
% $ Pfad: LaTeXTemplate/Handbuch Magic Faucet Fountain/Chapters/de/Hauptfunktion.tex $
% $ Kurzbeschreibung: Dieses Dokument zeigt die Hauptfunktion der Applikation. $
%
% !TeX spellcheck = de_DE/GB
% !TeX program = pdflatex
% !BIB program = biber/bibtex
% !TeX encoding = utf8
%
%%%%%%%%%%%%%%%%%%%%%%%%%%%%%%%%%%%%%%%%

\chapter{Hauptfunktion}
\begin{figure}[!ht]
	\centering
	\resizebox{0.5\textwidth}{!}{%
		\begin{circuitikz}
			\tikzstyle{every node}=[font=\LARGE]
			\draw [line width=1.2pt, short] (3.5,13.25) -- (6,3.25);
			\draw [line width=1.2pt, short] (22.5,13.25) -- (20,3.25);
			\draw [line width=1.2pt, short] (6,3.25) .. controls (9.75,0.75) and (16.25,0.75) .. (20,3.25);
			\draw [line width=1.2pt, short] (12.5,12.25) -- (12.5,26.75);
			\draw [line width=1.2pt, short] (13.5,12.25) -- (13.5,26.75);
			\draw [line width=1.2pt, short] (6.25,12) .. controls (6,13.25) and (10.75,13.5) .. (12.5,13.5);
			\draw [line width=1.2pt, short] (13.5,13.5) .. controls (17.75,13.5) and (20.25,12.75) .. (19.75,12);
			\draw [line width=1.2pt, short] (12.5,15.25) .. controls (0.5,14.75) and (2.25,11.25) .. (13,11.25);
			\draw [line width=1.2pt, short] (13.5,15.25) .. controls (24.25,15) and (24.75,11.25) .. (13,11.25);
			\draw [line width=1.2pt, short] (12.5,15.75) .. controls (2.25,15.5) and (-1.25,11.25) .. (13,10.75);
			\draw [line width=1.2pt, short] (13.5,15.75) .. controls (26.5,15) and (24.5,10.75) .. (13,10.75);
			\draw [line width=1.2pt, short] (12.5,12.25) .. controls (12.75,12) and (13.25,12) .. (13.5,12.25);
			\draw [ line width=1.2pt ] (19.5,7.25) circle (0.25cm);
			\draw [line width=1.2pt, short] (12,27) -- (12,27.5);
			\draw [line width=1.2pt, short] (14,27) -- (14,27.5);
			\draw [line width=1.2pt, short] (12.5,27.5) .. controls (12.25,30) and (13,30.5) .. (14.5,31.25);
			\draw [line width=1.2pt, short] (13.5,27.5) .. controls (13.25,29) and (13.5,29.5) .. (14.5,30);
			\draw [line width=1.2pt, short] (14.5,31.25) .. controls (15,31.25) and (15.25,31.5) .. (15.25,32);
			\draw [line width=1.2pt, short] (17.25,32) -- (17.25,32.5);
			\draw [line width=1.2pt, short] (16,32.5) -- (16,33.5);
			\draw [line width=1.2pt, short] (16.5,32.5) -- (16.5,33.5);
			\draw [ line width=1.2pt ] (16.25,34.75) ellipse (0.5cm and 0.25cm);
			\draw [line width=1.2pt, short] (15.75,34.75) -- (14.25,34.5);
			\draw [line width=1.2pt, short] (14.5,34) -- (16,34.25);
			\draw [line width=1.2pt, short] (15.75,34.25) -- (15.75,33.75);
			\draw [line width=1.2pt, short] (16.75,34.75) -- (16.75,33.75);
			\draw [line width=1.2pt, short] (16.75,34.5) -- (18,34.75);
			\draw [line width=1.2pt, short] (16.5,35) -- (17.75,35.25);
			\draw [line width=1.2pt, short] (14.25,34.5) .. controls (14.5,34.25) and (14.75,34.25) .. (14.5,34);
			\draw [line width=1.2pt, short] (14.25,34.5) .. controls (14,34) and (14.25,34) .. (14.5,34);
			\draw [line width=1.2pt, short] (17.75,35.25) .. controls (18,35) and (18.25,35) .. (18,34.75);
			\draw [line width=1.2pt, short] (15.75,33.75) .. controls (16,33.25) and (16.75,33.5) .. (16.75,33.75);
			\draw [line width=1.2pt, short] (16,32.75) .. controls (15,32.5) and (15,32.25) .. (16.25,32.25);
			\draw [line width=1.2pt, short] (16.5,32.75) .. controls (17.5,32.5) and (17.5,32.25) .. (16.25,32.25);
			\draw [line width=1.2pt, short] (15.25,32.5) -- (15.25,32);
			\draw [line width=1.2pt, short] (15.25,32) .. controls (15.75,31.5) and (17,31.75) .. (17.25,32);
			\draw [line width=1.2pt, short] (16,32.5) .. controls (16,32.25) and (16.5,32.25) .. (16.5,32.5);
			\draw [line width=1.2pt, short] (15,29.75) .. controls (16,29.25) and (16.25,29) .. (17.5,29.5);
			\draw [line width=1.2pt, short] (17.25,32) .. controls (18.25,31.5) and (18.75,30) .. (17.5,29.5);
			\draw [line width=1.2pt, short] (14.5,30) .. controls (14.5,30) and (14.75,30) .. (15,29.75);
			\draw [line width=1.2pt, short] (17.75,31.75) -- (19,32.25);
			\draw [line width=1.2pt, short] (18.25,30.75) -- (19.5,31.25);
			\draw [line width=1.2pt, short] (19,32.25) .. controls (19.25,32.25) and (19.75,31.25) .. (19.5,31.25);
			\draw [line width=1.2pt, short] (19,32.75) -- (19.75,33);
			\draw [line width=1.2pt, short] (19.75,30.75) -- (20.5,31);
			\draw [line width=1.2pt, short] (19,32.75) .. controls (20,32.75) and (20.5,30.75) .. (19.75,30.75);
			\draw [line width=1.2pt, short] (19,31) .. controls (19.25,30.75) and (19.5,30.5) .. (19.75,30.75);
			\draw [line width=1.2pt, short] (19,32.75) .. controls (19,32.5) and (18.75,32.5) .. (18.75,32.25);
			\draw [line width=1.2pt, short] (19.75,33) .. controls (20.5,32.75) and (21,31.5) .. (20.5,31);
			\draw [line width=1.2pt, short] (12.5,27.75) .. controls (11.25,27.25) and (12.5,27) .. (13,27.25);
			\draw [line width=1.2pt, short] (13.5,27.75) .. controls (14.75,27.25) and (13.5,27) .. (13,27.25);
			\draw [line width=1.2pt, short] (12,27) .. controls (12.25,26.75) and (12.5,26.75) .. (13,26.75);
			\draw [line width=1.2pt, short] (13,26.75) .. controls (13.5,26.75) and (13.75,26.75) .. (14,27);
			\draw [line width=1.2pt, short] (12.5,27.5) .. controls (12.75,27.25) and (13.25,27.25) .. (13.5,27.5);
		\end{circuitikz}
	}%
	\caption{Skizze des Magic Faucet Fountain}
	\label{fig:magicfaucet}
\end{figure}

Der Magic Faucet Fountain erzeugt die Illusion, wie in (Skizze) zu sehen, in der ein Wasserhahn durch eine Wassersäule in der Luft gehalten wird. Der Wasserhahn schwebt auf dem Wasser mithilfe eines Acryl-Rohres und begießt die dekorativen Steine, wie in (Skizze) zu sehen ist. Im Blumentopf befindet sich ein Wassertank, aus dem mithilfe einer Pumpe das Wasser gefördert wird. Der Wasserstand ist mit einer App auf dem Smartphone einsehbar. Zudem lässt sich die Pumpe ebenfalls durch die App ein- und ausschalten. Der Netzteilkasten ist zudem mit einem Hauptschalter ausgestattet, der das System ein- und ausschaltet. 

\section{Erste Schritte}
Sobald Sie Ihren Magic Faucet Fountain erhalten oder gekauft haben, können Sie ihn schon direkt nutzen! Packen Sie das Produkt einfach aus, schauen Sie nach ob alle Teile vorhanden sind (siehe Stückliste:\ref{stückliste}) und halten Sie Ihr Smartphone bereit. Sie benötigen kein extra Werkzeug oder unnötig komplizierte Anleitungen!

\section{Aufbau des Magic Faucet Fountain}
Nachdem Sie den Blumentopf und den Deckel ausgepackt haben, überprüfen Sie bitte, ob alle Bestandteile des Produkts vorhanden sind. Um den Magic Faucet Fountain aufzubauen, müssen Sie lediglich den Deckel auf den Blumentopf montieren und die dekorativen Steine auf den Deckel legen. Danach füllen Sie das Wasser bis zur Markierung auf. Abschließend können Sie das Netzkabel anschließen und mit dem Einrichten der App beginnen! Falls Sie genauere Informationen zum Aufbau benötigen, schauen Sie einfach in der Skizze \ref{fig:Skizze} nach!

\section{Einrichten der App}
Laden Sie sich die RemoteXY - App über den Google Playstore oder den Appstore auf Ihr Smartphone herunter. Schalten Sie die Bluetooth-Verbindung auf Ihrem Smartphone ein und suchen Sie nach in der Nähe befindlichen Geräten. Schalten Sie den Netzwerkkasten über den Kippschalter ein und überprüfen Sie, ob die grüne LED leuchtet. Ist dies der Fall warten Sie auf das Blinken der blauen LED. Sobald dies der Fall ist können Sie in der RemoteXY-App die Verbindung mit dem Magic Faucet Fountain herstellen. Die Verbindung ist hergestellt, sobald die blaue LED konstant leuchtet. Nun können Sie Ihren Magic Faucet Fountain nutzen!
Die App, wie in \ref{fig:ansichtApp} dargestellt, verfügt über eine Darstellung des Blumentopfes und eine Füllstandsanzeige. Zudem lässt sich die Pumpe mit zwei Knöpfen ansteuern.


\begin{figure}[]
		\centering
		\resizebox{1\textwidth}{!}{%
			\begin{circuitikz}
				\tikzstyle{every node}=[font=\LARGE]
				
				% Gehäuse
				\draw [fill={rgb,255:red,128; green,128; blue,128}, line width=2pt] (-5,12) rectangle (12.5,7);
				\draw [fill={rgb,255:red,128; green,128; blue,128}, line width=1pt] (-5,12) -- (12.5,12) -- (15,14.5) -- (-2.5,14.5) -- cycle;
				\draw [line width=2pt](-5,12) -- (-2.5,14.5);
				\draw [line width=2pt](-2.5,14.5) -- (15,14.5);
				\draw [line width=2pt](12.5,12) -- (15,14.5);
				\draw [line width=2pt](12.5,7) -- (15,9.5);
				\draw [line width=2pt](15,9.5) -- (15,14.5);
				\draw [fill={rgb,255:red,128; green,128; blue,128}, line width=1pt] (12.5,7) -- (15,9.5) -- (15,14.5) -- (12.5,12) -- cycle;
				
				% Obere LED-Kappen (Ellipse oben)
				\draw [color=blue, fill=blue, line width=2pt] (0,13.25) ellipse (1.25cm and 0.25cm);
				\draw [color=green!50!black, fill=green, line width=2pt] (3.75,13.25) ellipse (1.25cm and 0.25cm);
				\draw [fill=gray!50, line width=2pt] (8.75,13.25) ellipse (1.25cm and 0.25cm);
				
				% Zylinderkörper der LEDs
				\draw [color=blue, fill=blue] (-1.25,13.25) rectangle (1.25,15);
				\draw [color=green!50!black, fill=green] (2.5,13.25) rectangle (5,15);
				% Beschriftungen über den Komponenten
				\node at (0,17.5) {\textbf{LED Blau}};
				\node at (3.75,17.5) {\textbf{LED Grün}};
				\node at (8.75,17.5) {\textbf{Kippschalter}};
				\draw [->, >=Stealth] (0,17) -- (0,15.5);
				\draw [->, >=Stealth] (3.75,17) -- (3.75,15.5);
				\draw [->, >=Stealth] (9.25,17) -- (9.25,16);
				% Untere Halbellipsen (schwarze Kurven)
				\draw [color=black, line width=1pt, domain=-1.25:1.25, samples=100] 
				plot (\x, {-0.25 * sqrt(1 - (\x/1.25)^2) + 13.25});
				\draw [color=black, line width=1pt, domain=2.5:5, samples=100] 
				plot (\x, {-0.25 * sqrt(1 - ((\x - 3.75)/1.25)^2) + 13.25});
				
				% Oberer Ellipsenrand (Top-Kappe)
				\draw [fill=green, line width=2pt] (3.75,15) ellipse (1.25cm and 0.25cm);
				\draw [fill=blue, line width=2pt] (0,15) ellipse (1.25cm and 0.25cm);
				
				% Kippschalter
				\draw [fill=gray!50, line width=0.8pt] (8.25,13.25) -- (9.25,13.25) -- (9.75,15.5) -- (8.75,15.5) -- cycle;
				\draw [fill=gray!50, line width=0.5pt] (9.25,15.5) ellipse (0.5cm and 0.25cm);
				%Kabel
				\draw [line width=10pt, short] (13.75,9.5) .. controls (19,10.5) and (17.75,8.5) .. (21.25,9.5);
				\draw [line width=10pt, short] (-5,9.5) .. controls (-7.5,10.75) and (-7.75,8.25) .. (-10,9.5);
			\end{circuitikz}
		}%
	
		\caption{Netzteil des Magic Faucet Fountain}
	\label{fig:ansichtApp}
\end{figure}
	


\begin{figure}[]
	\centering
	\resizebox{0.75\textwidth}{!}{%
		\begin{circuitikz}
			\tikzstyle{every node}=[font=\LARGE]
			\draw [ fill={rgb,255:red,224; green,234; blue,237} , line width=1.3pt , rounded corners = 41.4] (-12.5,32) rectangle (2.5,9.5);
			\draw [line width=1.3pt, short] (-12.5,20.75) -- (2.5,20.75);
			\draw [ fill={rgb,255:red,122; green,122; blue,122} , line width=1.3pt ] (-11.25,19.5) rectangle (1.25,17);
			\draw [line width=1.3pt, short] (-10,29.5) -- (0,29.5);
			\draw [line width=1.3pt, short] (-10,29.5) -- (-10,28.25);
			\draw [line width=1.3pt, short] (-10,28.25) -- (-8.75,28.25);
			\draw [line width=1.3pt, short] (0,29.5) -- (0,28.25);
			\draw [line width=1.3pt, short] (0,28.25) -- (-1.25,28.25);
			\draw [line width=1.3pt, short] (-8.75,28.25) -- (-7.5,22);
			\draw [line width=1.3pt, short] (-7.5,22) -- (-2.5,22);
			\draw [line width=1.3pt, short] (-2.5,22) -- (-1.25,28.25);
			\draw [line width=1.3pt, short] (-8.25,29.5) -- (-7,22.75);
			\draw [line width=1.3pt, short] (-7,22.75) -- (-3,22.75);
			\draw [line width=1.3pt, short] (-3,22.75) -- (-1.5,29.5);
			\draw [ fill={rgb,255:red,166; green,187; blue,191} , line width=1.3pt , rounded corners = 30.0] (-11.25,14.5) rectangle  node {\Huge Start} (-6.25,10.5);
			\draw [ fill={rgb,255:red,166; green,187; blue,191} , line width=1.3pt , rounded corners = 30.0] (-3.75,14.5) rectangle  node {\Huge Stop} (1.25,10.5);
			\begin{scope}[rotate around={-2.5:(-7.75,26.75)}]
				\draw[domain=-7.75:-2.25,samples=100,color={rgb,255:red,36; green,24; blue,248}, , line width=2pt] plot (\x,{0.3*sin(1.01*\x r +7.55 r ) +26.75});
			\end{scope}
			\draw [line width=1.3pt, short] (-1.25,27) -- (1.25,27)node[pos=0.5, fill=white]{Max};
			\draw [line width=1.3pt, short] (-1.25,23.25) -- (1.25,23.25)node[pos=0.5, fill=white]{Min};
			\node [font=\LARGE] at (-5,24.5) {200 ml};
			\draw [ fill={rgb,255:red,249; green,224; blue,49} , line width=1.3pt , rounded corners = 9.0] (-11,19.25) rectangle (-10.25,17.25);
			\draw [ fill={rgb,255:red,249; green,224; blue,49} , line width=1.3pt , rounded corners = 9.0] (-9.75,19.25) rectangle (-9,17.25);
			\draw [ fill={rgb,255:red,249; green,224; blue,49} , line width=1.3pt , rounded corners = 9.0] (-8.5,19.25) rectangle (-7.75,17.25);
			\draw [ fill={rgb,255:red,249; green,224; blue,49} , line width=1.3pt , rounded corners = 9.0] (-7.25,19.25) rectangle (-6.5,17.25);
			\draw [ fill={rgb,255:red,249; green,224; blue,49} , line width=1.3pt , rounded corners = 9.0, rotate around={-360:(-5.625, 18.25)}] (-6,19.25) rectangle (-5.25,17.25);
			\draw [ fill={rgb,255:red,122; green,122; blue,122} , line width=1.3pt ] (1.25,18.75) rectangle (1.75,17.75);
			\draw [line width=1pt, ->, >=Stealth] (6.5,27.25) -- (-0.25,24.75);
			\draw [line width=1pt, ->, >=Stealth] (6.25,21.75) -- (1.25,19.5);
			\draw [line width=1pt, ->, >=Stealth] (-5.75,6.75) -- (-6.75,10.5);
			\draw [line width=1pt, ->, >=Stealth] (-4.25,6.75) -- (-3.5,10.5);
			\node [font=\LARGE] at (10.75,27.75) {Genaue Wasserstandanzeige};
			\node [font=\LARGE] at (11,22.25) {Farbliche Wasserstandanzeige};
			\node [font=\LARGE] at (-5.5,5) {Magic Faucet Fountain};
			\node [font=\LARGE] at (-5.25,6) {Tasten zum Ein- und Ausschalten des };
			\draw [ fill={rgb,255:red,123; green,230; blue,87} , line width=1pt , rounded corners = 9.0] (6.25,20.75) rectangle (7.5,19.5);
			\draw [ fill={rgb,255:red,249; green,224; blue,49} , line width=1pt , rounded corners = 9.0] (6.25,18.25) rectangle (7.5,17);
			\draw [ fill={rgb,255:red,240; green,34; blue,23} , line width=1pt , rounded corners = 9.0] (6.25,15.75) rectangle (7.5,14.5);
			\node [font=\LARGE] at (14.5,17.5) {Wasser sinkt. Bei Gelegenheit nachfüllen.};
			\node [font=\LARGE] at (14.75,15) {Zu wenig Wasser. Bitte Wasser nachfüllen!};
			\node [font=\LARGE] at (15,20) {Wasserstand normal. Keine Maßnahmen nötig.};
		\end{circuitikz}
	}%
	\caption{App-Ansicht des Füllstands und Pumpensteuerung}
    \label{fig:appAnsicht}
\end{figure}

